% Part: sets-functions-relations
% Chapter: functions
% Section: isomorphisms

\documentclass[../../../include/open-logic-section]{subfiles}

\begin{document}

\olfileid{sfr}{fun}{iso}
\olsection{એકરૂપતા}

\begin{explain}
\emph{એકરૂપતા} એ જોડાતા ગણોની સંરચના જાળવતું એક-એક વ્યાપ્ત વિધેય છે.
સંરચના એટલે ગણોના !!{element}s વચ્ચે રહેલા સંબંધો.
બે ગણ $X=\{1,2,3\}$ અને $Y=\{4,5,6\}$ લો.
બંનેમાં અનુગામી હોવાનો, નાનું હોવાનો અને મોટું હોવાનો સંબંધ છે.
બે ગણો વચ્ચેની એકરૂપતા આ સંરચનાઓ જાળવતું !!a{bijection} છે.
એટલે !!a{bijective} ફલન~$f \colon X
\to Y$ એકરૂપતા છે જો, ગણના કોઈ પણ બે ઘટકો માટે,
$i<j$ જો અને ફક્ત જો $f(i)<f(j)$,
$i>j$ જો અને ફક્ત જો $f(i)>f(j)$,
અને $j$ એ $i$નો અનુગામી હોય જો અને ફક્ત જો $f(j)$ એ
$f(i)$નો અનુગામી હોય.
\end{explain}

\begin{defn}[એકરૂપતા]
$U$ એ જોડી $\langle X, R\rangle$ અને $V$ એ જોડી
$\langle Y, S\rangle$ હોય, જ્યાં $X$ અને $Y$ ગણો છે અને
$R$ તથા $S$ અનુક્રમે $X$ તથા $Y$ પરના સંબંધો છે.
$X$થી $Y$માં જતું !!{bijection}~$f$ એ $U$થી $V$માં જતી
\emph{એકરૂપતા} છે જો અને ફક્ત જો તે સંબંધની સંરચના જાળવે:
$X$ના કોઈ પણ $x_{1}$ અને $x_{2}$ માટે
$\tuple{x_1,x_2} \in R$ જો અને ફક્ત જો
$\tuple{f(x_1),f(x_2)} \in S$.
\end{defn}

\begin{ex}
બે ગણ $X=\{1,2,3\}$ અને $Y=\{4,5,6\}$ લો,
સાથે નાનું હોવાનો અને મોટું હોવાનો સંબંધ લો.
$f(x) = 7-x$ વડે આપેલું ફલન~$f\colon X \to Y$
એ $\tuple{X,<}$ અને $\tuple{Y,>}$ વચ્ચેની એકરૂપતા છે.
\end{ex}

\end{document}

